% =====================================================================
% sections/related_work.tex
% =====================================================================
\section{Background and Related Work}\label{sec:related}

\PG{} sits where three lines of work meet: mobile GUI agents and the attacks against them, system-level defenses for LLM agents, and operating-system mechanisms for information flow and UI integrity. We review each, and close with a comparison (\cref{tab:related}) that fixes the boundary of our contribution.

\subsection{Mobile GUI agents and how they perceive the screen}\label{sec:rw-agents}

Early LLM-driven phone agents read a textual rendering of the view hierarchy. AutoDroid converts the GUI into a simplified HTML representation derived from the view hierarchy and reaches a 71.3\% task success rate on 158 common tasks~\cite{wen2024autodroid}; AutoDroid-V2, presented at MobiSys 2025, moves to on-device small language models that generate scripts over an app-documentation abstraction~\cite{wen2025autodroidv2}. AppAgent learns app usage by exploration and acts through a simplified action space over annotated screenshots~\cite{zhang2025appagent}. A second line of work removes system metadata from the loop. Mobile-Agent is vision-centric and avoids system-specific customization~\cite{wang2024mobileagent}; Mobile-Agent-v2 adds planning, decision, and reflection agents~\cite{wang2024mobileagentv2}; and Mobile-Agent-v3 builds a multi-agent framework on the GUI-Owl foundation model~\cite{ye2025mobileagentv3}. CogAgent~\cite{hong2024cogagent}, SeeClick~\cite{cheng2024seeclick}, and OS-ATLAS~\cite{wu2025osatlas} train vision-language models for GUI grounding from screenshots, and UI-TARS is a native agent model that ``solely perceives the screenshots as input''~\cite{qin2025uitars}. AutoGLM decouples planning from grounding through an intermediate interface~\cite{liu2024autoglm}, and its open release Open-AutoGLM operates from screenshots over ADB~\cite{openautoglm2025software}. DroidRun, recently renamed Mobilerun, combines accessibility trees with screenshots and is model-agnostic~\cite{droidrun2025software}. At MobiSys 2026, AgentProg manages a GUI agent's context with program-guided state~\cite{tian2026agentprog}. Surveys organize this space by environment, perception, and action execution~\cite{zhang2025llmbrainedguiagents}, and a separate survey treats trustworthiness, including security, as a first-class dimension~\cite{shi2025trustworthyguiagents}.

Benchmarks have followed. Android in the Wild provides 715k episodes of device control~\cite{rawles2023aitw}. AndroidWorld offers 116 programmatic tasks across 20 real apps with state-based success checks on a live emulator~\cite{rawles2025androidworld}. AndroidLab contributes 138 tasks across nine apps and supports both accessibility-tree (XML) and screenshot (set-of-mark) modes~\cite{xu2025androidlab}, and SPA-Bench integrates more than ten agents with metrics that include resource consumption~\cite{chen2025spabench}. We use AndroidWorld and AndroidLab to measure the utility cost of \PG{}.

Two observations from this literature shape our design. First, perception modality varies: some agents read the accessibility tree, some read only pixels, and some read both. Provenance therefore has to be attached to \emph{both} channels. Second, every agent above treats the screen as a flat input; none distinguishes regions by the principal that produced them.

\subsection{Indirect prompt injection against GUI agents}\label{sec:rw-attacks}

Greshake et al.\ showed that LLM-integrated applications ``blur the line between data and instructions'' and introduced indirect prompt injection through retrieved content~\cite{greshake2023notwhat}. Benchmarks for tool-calling agents followed: InjecAgent with 1,054 test cases~\cite{zhan2024injecagent} and AgentDojo with 97 tasks and 629 security test cases~\cite{debenedetti2024agentdojo}. Liu et al.\ formalized prompt injection as modifying an application's data so that it accomplishes an injected task, and found that no existing prevention-based defense is sufficient~\cite{liu2024formalizing}.

For agents that read the screen, the environment itself is the injection channel. On the web, EIA injects malicious elements adapted to blend into the page and steals specific personal information with up to 70\% success~\cite{liao2025eia}; adversarial pop-ups reach an average attack success rate of 86\% on OSWorld and VisualWebArena while basic defenses such as instructing the agent to ignore pop-ups are ineffective~\cite{zhang2025popups}; and even benign but distracting content misleads multimodal agents~\cite{ma2025caution}. WASP measures end-to-end web-agent security against low-effort injections~\cite{evtimov2025wasp}, Wu et al.\ show that a trigger image covering less than 5\% of a page's pixels can redirect multimodal agents~\cite{wu2025dissecting}, and Chen et al.\ show that commands buried in dense, legalistic text are parsed indiscriminately by web agents~\cite{chen2025fineprint}.

On phones, the evidence is now extensive and recent. AgentHazard (MobiSys 2026) instruments real apps so that only regions writable by third parties are modified, and reports average misleading rates of 42.0\% (dynamic, 122 tasks) and 36.1\% (static, more than 3,000 scenarios)~\cite{liu2026agenthazard}. Du et al.\ measure eight agents against injections through in-app ads, WebViews, e-mail, and messaging notifications, and find that interstitial ads lead to full execution of the malicious workflow in 93.75\% of trials~\cite{du2025measuringmobileagents}. GhostEI-Bench contributes 110 test cases of dynamic and static environmental injection, including overlays and pop-up SMS, with vulnerability rates between 16.43\% and 56.25\%~\cite{chen2025ghosteibench}. MobileWorldSafety covers five injection vectors embedded in app workflows across 13 apps~\cite{chen2026mobileworldsafety}. Deivasigamani et al.\ document nine accessibility-exposed vectors, A1 to A9, including notification injection, cross-application overlays, hidden accessibility metadata, and HTML injection~\cite{deivasigamani2026nota11y}. Li et al.\ define human--agent UI desynchronization, a cross-observer inconsistency in which content is ``available or salient to an agent but unavailable, inconspicuous, or differently represented to the user,'' and realize it through display-cutout occlusion, low luminance contrast, accessibility-only content, and impractically small or transparent widgets~\cite{li2026uidesync}. Qian et al.\ add a temporal dimension: action rebinding exploits the observe-to-act gap (4.18\,s to 15.43\,s in their measurements) by switching the foreground with a standard Intent, and they trace the root cause to the fact that GUI agents ``lack a mechanism to bind their actions to a specific process handle or a unique component identity''~\cite{qian2026zeropermission}.

Two adjacent desktop works are close to parts of \PG{}. PUSV re-verifies the UI state immediately before each dispatched action through masked pixel similarity at the click target and screenshot and X-window snapshot differences, intercepting notification-overlay and focus-manipulation attacks, but it has a structural blind spot for injections with no visual footprint~\cite{xu2026pusv}. The visual confused deputy work defines the event in which an agent's click lands on an element other than the one it intended, and guards clicks with a learned dual-channel classifier outside the agent process~\cite{liu2026visualconfuseddeputy}. Both are desktop systems, neither carries origin labels, and the second places a learned model in the decision path. \PG{}'s \DB{} is closest to PUSV in spirit but binds to OS identities (window token, node identity, node digest) rather than to pixels.

\subsection{Defenses for LLM agents}\label{sec:rw-defenses}

\paragraph{Prompt-level and model-level defenses}
Spotlighting transforms untrusted input through delimiting, datamarking, or encoding to give the model ``a reliable and continuous signal of its provenance''~\cite{hines2024spotlighting}. The signal is advisory: the model may ignore it, and delimiters are among the prevention defenses that Liu et al.\ found insufficient~\cite{liu2024formalizing}. StruQ separates instruction and data channels with a secure front-end and a fine-tuned model~\cite{chen2025struq}, SecAlign applies preference optimization against injected inputs~\cite{chen2025secalign}, and the instruction hierarchy trains models to prioritize privileged instructions over tool outputs and over instructions in images~\cite{wallace2024instruction}. All three need training access and remain probabilistic. In our threat model the agent's model is untrusted, so these defenses are complements that \PG{} assumes may fail.

\paragraph{Detection}
MELON detects indirect prompt injection in tool-calling agents by re-executing each step with the user task replaced by a task-neutral prompt and flagging an attack when the tool calls of the original and masked runs are similar~\cite{zhu2025melon}. It doubles model calls and is evaluated on AgentDojo with temperature zero. Its ``provable'' guarantee is a Hoeffding-type error bound for an ensemble of masking prompts under a separation assumption~\cite{zhu2025melon}. \PG{} borrows the masked re-execution idea but inverts it: it masks \emph{untrusted screen regions selected by OS labels} and treats a \emph{change} of the protected action as evidence of untrusted causation. OS-Sentinel combines a formal verifier over system-state traces with a VLM judge over observation--action pairs, at the step level and the trajectory level~\cite{sun2026ossentinel}; it classifies actions as unsafe but does not attribute them to the content that caused them, and it keeps a VLM judge in the decision path. NeuroTaint tracks taint through LLM agents with a provenance graph, semantic matching, and an LLM-based causal analyzer~\cite{cai2026neurotaint}; it too relies on an LLM judge.

\paragraph{System-level designs for tool-calling agents}
IsolateGPT isolates third-party apps in spokes mediated by a hub with deterministic, non-LLM operators~\cite{wu2025isolategpt}. ACE splits planning into an abstract plan derived only from the trusted user query and a concrete plan checked against it and against a lattice-based information-flow policy~\cite{li2026ace}. CaMeL extracts control and data flow from the trusted query and tracks capabilities on every value, so untrusted data ``can never impact the program flow''~\cite{debenedetti2025camel}. Fides attaches confidentiality and integrity labels to all data an agent processes and provides a non-interference guarantee for integrity of consequential actions~\cite{costa2025fides}. Permissive IFC propagates only the labels of inputs that were influential for an output, testing influence by substituting the context~\cite{siddiqui2025permissive}, an idea related to our masked re-execution but applied to confidentiality labels of text outputs. Progent enforces least-privilege policies over tool names and arguments~\cite{shi2025progent}, APPA branches contexts to confine taint while preserving utility~\cite{kravchenko2026appa}, and FlowSeal enforces confidentiality outside the model's context with a provenance-derived lattice, arguing that ``whenever enforcement is a judgment the LLM makes over the same conversational context an adversary controls, the enforcement mechanism and the attack surface coincide''~\cite{shim2026flowseal}. We share that principle. These systems, however, operate on discrete tool outputs whose source is known by construction, and several of them fix control flow before untrusted data arrives. Neither condition holds for a GUI agent.

\paragraph{Mobile agent defenses}
AutoDroid elides private information from prompts and asks the LLM to identify risky actions that require user confirmation~\cite{wen2024autodroid}, an in-model judgment. Aura proposes an intent-centric agent operating system with \texttt{TAG\_VERIFIED} and \texttt{TAG\_TAINTED} labels and a no-write-down policy into critical operations, evaluated on a simulated device in which app agents interact through structured mock APIs~\cite{zou2026aura}. AOHP implements an OS-level agent harness on AOSP with typed placeholders and taint tracking for sensitive data~\cite{zhao2026aohp}; its labels protect confidentiality and it does not evaluate injection. \PG{} is complementary to both: it keeps the GUI as the action surface, computes labels from real OS state, and enforces integrity of protected actions.

\subsection{Information flow and provenance in operating systems}\label{sec:rw-ifc}

Denning's lattice model gives information-flow policies their algebraic form~\cite{denning1976lattice}, and Biba's integrity model introduces the dual, integrity-oriented ordering~\cite{biba1977integrity}. Goguen and Meseguer define non-interference~\cite{goguen1982noninterference}, and Sabelfeld and Myers survey language-based enforcement~\cite{sabelfeld2003language}. Decentralized labels~\cite{myers1997decentralized} and Flume's OS-level decentralized IFC~\cite{krohn2007flume} showed that labels can be retrofitted onto standard OS abstractions. On Android, TaintDroid tracks privacy-sensitive data through variables, methods, messages, and files with modest overhead~\cite{enck2010taintdroid}, and CamFlow captures whole-system provenance in the kernel~\cite{pasquier2017camflow}. \PG{} borrows the lattice and the integrity orientation, but applies labels to \emph{perceptual} objects (accessibility nodes and screenshot tiles), and enforces them only at protected sinks through a causal test rather than through full propagation, which an opaque model makes impossible.

\subsection{Android UI security}\label{sec:rw-ui}

Accessibility has been an attack surface for over a decade. Jang et al.\ show that OS widgets handle real and accessibility input through the same channels, making the two hard to distinguish~\cite{jang2014a11y}. Cloak and Dagger combines \texttt{SYSTEM\_ALERT\_WINDOW} with an accessibility service to gain complete control of the UI feedback loop~\cite{fratantonio2017cloak}. Bianchi et al.\ study UI deception and propose an indicator of the app the user is interacting with~\cite{bianchi2015whattheapp}, and Fernandes et al.\ revisit those attacks and defenses~\cite{fernandes2016uideception}. Yan et al.\ analyze overlay malware at market scale at MobiSys 2019 and note that ``a programmatically visible overlay can be visually invisible to users''~\cite{yan2019overlay}, the gap our perceivability bit measures. Kalysch et al.~\cite{kalysch2018a11y} and AccessiLeaks~\cite{naseri2019accessileaks} document accessibility misuse for reading other apps' content. Classic web clickjacking, whose defense introduced the notions of visual and temporal context integrity for a user's sensitive action~\cite{huang2012clickjacking}, permission re-delegation~\cite{felt2011redelegation}, the confused deputy~\cite{hardy1988confused}, and time-of-check-to-time-of-use races~\cite{bishop1996toctou} are the conceptual ancestors of the attacks on agents: an agent with the user's authority is a deputy that untrusted content tries to confuse, and the observe-to-act gap is a check-to-use window.

Android has responded with point mechanisms that \PG{} generalizes. Since Android 12, the system blocks touches that pass through an untrusted overlay whose combined opacity exceeds a threshold (0.8 by default), while accessibility, input-method, and assistant windows are trusted~\cite{android12_untrusted_touch}. Android 15 prevents untrusted notification listeners from reading unredacted content of notifications that contain a detected one-time password~\cite{android15_otp_redaction}. Views can be marked as accessibility-data sensitive, in which case only services that declare \texttt{isAccessibilityTool} may read them~\cite{android_blog_a11y_sensitive,android_accessibilityserviceinfo}. AppFunctions lets apps expose functions to agents as tools, bypassing the GUI where apps opt in~\cite{android_appfunctions}. Each of these is a per-consumer or per-touch decision; none labels what an agent perceives or reasons about what an agent does with it.

\subsection{Causal attribution}\label{sec:rw-causal}

The masked re-execution test is an intervention in the sense of structural causal models: it sets the content of untrusted regions to a neutral value and observes whether the decision changes, a but-for test~\cite{pearl2009causality}. Cheaper attribution methods exist: integrated gradients satisfy sensitivity and implementation invariance but require gradient access~\cite{sundararajan2017axiomatic}, and attention rollout and attention flow approximate token importance from attention weights~\cite{abnar2020attentionflow}. Attention weights are not, in general, faithful explanations of a model's output~\cite{jain2019attention}, which is why \PG{} uses attention only as an optional pre-filter and never as the basis of its guarantee.

\subsection{Positioning}\label{sec:rw-position}

\Cref{tab:related} summarizes the boundary. To the best of our knowledge, within the literature verified for this manuscript (\texttt{citation\_search\_log.md}), no prior system computes origin and perceivability labels for accessibility nodes and screenshot regions in the OS, tests causal dependence of GUI actions on untrusted regions without an LLM judge, and binds actions to observed OS identities at dispatch. This statement is limited to the sources we verified; a fresh collision check against concurrent work is listed in \texttt{TODO\_before\_submission.md}.

\begin{table}[t]
\centering
\caption{Positioning against the closest prior work. \checkmark: provided; $\circ$: partial or advisory; --: not provided. ``Real Android'' means an implementation on an Android build or device. ``Labels inputs'' means that untrusted inputs to the agent carry provenance labels. ``No LLM in TCB'' means that no LLM/VLM judgment is part of the enforcement decision.}
\label{tab:related}
\footnotesize
\setlength{\tabcolsep}{3.5pt}
\begin{tabular}{@{}l c c c c c c c@{}}
\toprule
System & \makecell{Real\\Android} & \makecell{GUI\\agent} & \makecell{Labels\\inputs} & \makecell{Perceiv-\\ability} & \makecell{Causal\\test} & \makecell{Dispatch\\binding} & \makecell{No LLM\\in TCB} \\
\midrule
Spotlighting~\cite{hines2024spotlighting}          & -- & -- & $\circ$ & -- & -- & -- & -- \\
MELON~\cite{zhu2025melon}                          & -- & -- & -- & -- & \checkmark & -- & \checkmark \\
CaMeL~\cite{debenedetti2025camel}                  & -- & -- & \checkmark & -- & -- & -- & \checkmark \\
Fides~\cite{costa2025fides}                        & -- & -- & \checkmark & -- & -- & -- & \checkmark \\
ACE~\cite{li2026ace}                               & -- & -- & \checkmark & -- & -- & -- & \checkmark \\
NeuroTaint~\cite{cai2026neurotaint}                & -- & -- & \checkmark & -- & $\circ$ & -- & -- \\
OS-Sentinel~\cite{sun2026ossentinel}               & \checkmark & \checkmark & -- & -- & -- & -- & -- \\
Aura~\cite{zou2026aura}                            & -- & $\circ$ & \checkmark & -- & -- & -- & \checkmark \\
AOHP~\cite{zhao2026aohp}                           & \checkmark & \checkmark & $\circ$ & -- & -- & -- & \checkmark \\
PUSV~\cite{xu2026pusv}                             & -- & \checkmark & -- & -- & -- & $\circ$ & \checkmark \\
Visual Confused Deputy~\cite{liu2026visualconfuseddeputy} & -- & \checkmark & -- & -- & -- & $\circ$ & -- \\
\midrule
\PG{} (this work)                                  & $\circ$\textsuperscript{a} & \checkmark & \checkmark & \checkmark & \checkmark & \checkmark & \checkmark \\
\bottomrule
\end{tabular}

\smallskip
\raggedright\scriptsize \textsuperscript{a}\,Designed for AOSP on Pixel devices; the implementation is planned and not yet complete (\cref{sec:impl-plan}).
\end{table}
